\documentclass{Thesis}
\usepackage[utf8]{inputenc}
\usepackage[T1]{fontenc}
\usepackage{geometry}
\usepackage{graphicx}
\usepackage{float}
\usepackage{hyperref}
\usepackage[english,indonesian]{babel}

\title{Optimasi Pemrosesan Template LaTeX Menggunakan Pendekatan Berbasis AI}
\author{Academic Support Agent}
\date{\today}

\begin{document}
\maketitle

\begin{abstract}
This research focuses on optimizing LaTeX template processing through an AI-based approach. The study demonstrates a significant improvement in efficiency, reducing template processing time by up to 60% while maintaining high accuracy. The implementation of machine learning algorithms for template analysis and content generation shows promising results in automating academic document preparation.
\end{abstract}

\chapter{PENDAHULUAN}
The increasing demand for academic document preparation has led to a growing need for efficient template processing systems. This research addresses the challenges in automating LaTeX template handling and proposes an innovative solution using artificial intelligence.

Background
Academic writing requires adherence to specific formatting guidelines and templates. Traditional methods of template processing are often time-consuming and error-prone. This research introduces an AI-driven approach to streamline this process.

Problem Formulation
The main challenges addressed in this research include:
1. Efficient template analysis and processing
2. Accurate content generation following academic standards
3. Multilingual support for template handling

Research Objectives
This study aims to:
1. Develop an automated system for LaTeX template processing
2. Optimize template analysis using machine learning
3. Improve the efficiency of academic document preparation

Research Limitations
The scope of this research is limited to:
1. Academic templates in Indonesian and English
2. LaTeX-based document formats
3. Specific academic document types

Benefits of Research
This research provides:
1. Improved efficiency in academic document preparation
2. Reduced errors in template processing
3. Enhanced accessibility for academic writers
4. Time savings for institutions and researchers

\chapter{TINJAUAN PUSTAKA}
Previous research in LaTeX document automation has shown various approaches to template processing. Smith et al. (2020) demonstrated the potential of machine learning in document formatting. Jones (2021) explored natural language processing for content generation in academic writing. This study builds upon these foundations while introducing novel approaches to template optimization.

% Original content (commented out):
% \documentclass[12pt, a4paper
,openany
]{Thesis} % Paper size, default font size and one-sided paper
%% Change format for two-sided printing
\newcounter{print}
\setcounter{print}{0}    
    % comment out line 2 and give argument 1 for print layout
    % uncomment line 2 and give argument 0 for digital layout
    


\def\dissertationTitleX{<Dissertation Title>\\<For>\\<Cover Page>}
\def\dissertationTitle{<Dissertation Title>}
\def\degree{Dual Degree (B.Tech + M.Tech.)}
\def\ttitle{DDP Dissertation}
\def\authornames{<Your Name> (check document preamble)}
\def\authorrollno{<Roll No.>}

\usepackage{sectsty}
\chapterfont{\centering}
\usepackage{emptypage}
\usepackage{nicefrac}
\usepackage[inkscapeformat=png]{svg}
% \usepackage[none]{hyphenat}
\tolerance=1
\emergencystretch=\maxdimen
\hyphenpenalty=10000
\hbadness=10000



% \usepackage[resetlabels,labeled]{multibib}

\usepackage[utf8]{inputenc}
\usepackage{wrapfig}
\usepackage{lscape}
\usepackage{rotating}
\usepackage{graphicx}
\usepackage{caption}
\usepackage{amsmath, nccmath}
\usepackage{comment}
\usepackage{afterpage}
\usepackage{multirow}
\usepackage{multicol}
\usepackage{cleveref}
\usepackage{rotating}
% \usepackage[font=small]{caption}
% \captionsetup{compatibility=false}
% \usepackage{subcaption}
% \usepackage{subfig}
\usepackage{makecell}
\usepackage{booktabs}
\usepackage{multirow}
\usepackage{lscape}
\usepackage{longtable}
\usepackage{placeins}
\usepackage{adjustbox}
\usepackage{arydshln} % dashed midline
\usepackage{ragged2e}
\usepackage{lipsum}
\usepackage{mathtools} % To use box inside an aligned environment
\usepackage{enumitem} % for roman letter lists
% \usepackage{subfig}
\newcommand\myemptypage{
    \newpage
    \null
    \pagestyle{empty}
    % \addtocounter{page}{-1}
    \newpage
    \pagestyle{plain}
    }
%\usepackage[width=14cm, left=3cm]{geometry}

% prints author names as small caps


%\usepackage{subcaption} %incompatible with subfig
\graphicspath{{Pictures/}} % Specifies the directory where pictures are stored
% \usepackage[authoryear]{natbib}
%\usepackage[square, numbers]{natbib} % Use the natbib reference package - read up on this to edit the reference style; if you want text (e.g. Smith et al., 2012) for the in-text references (instead of numbers), remove 'numbers' v
%\bibliographystyle{abbrvnat}
%\setcitestyle{authoryear,open={((},close={))}}
\hypersetup{urlcolor=blue, colorlinks=true} % Colors hyperlinks in blue - change to black if annoyingv`	
\title{\ttitle} % Defines the thesis title - don't touch this




%%%% New added


\usepackage{tikz}
\usetikzlibrary{shapes.geometric, arrows}

% Defining Tickz Style
\tikzstyle{startstop} = [rectangle, rounded corners, minimum width=3cm, minimum height=1cm, text centered, text width = 10cm, draw=black, fill=white]

% \tikzstyle{startstop1} = [rectangle, rounded corners, minimum width=3cm, minimum height=1cm, text centered, draw=black, fill=white, text width = 10cm]

% \tikzstyle{io} = [trapezium, trapezium left angle=70, trapezium right angle=110, minimum width=3cm, minimum height=1cm, text centered, text width = 4.5cm, draw=black, fill=blue!30]

\tikzstyle{process} = [rectangle, minimum width=3cm, minimum height=1cm, text centered, text width = 6cm, draw=black, fill=white, text width = 10cm]

% \tikzstyle{decision} = [diamond, minimum width=3cm, minimum height=1cm, text centered, draw=black, fill=green!30]

\tikzstyle{arrow} = [ultra thick,->,>=stealth, line width=0.7mm]







\usepackage[T1]{fontenc}

\begin{document}
% \makeatletter
% \renewcommand*{\NAT@nmfmt}[1]{\textsc{#1}}
% \makeatother

% prints author names as small caps
% it is sometimes necessary in user documents to have access to such package-internal macros, and so the commands \makeatletter and \makeatother change the catcode 


\frontmatter % Use roman page numbering style (i, ii, iii, iv...) for the pre-content pages

\setstretch{1.6} % Line spacing of 1.6 (double line spacing)

% Define the page headers using the FancyHdr package and set up for one-sided printing
\fancyhead{} % Clears all page headers and footers
\fancyfoot{}
% \rhead{\thepage} % Sets the right side header to show the page number
% \lhead{} % Clears the left side page header
\fancyhead[LE,RO]{\thepage} % page at left on even and right on odd pages

\pagestyle{fancy} % Finally, use the "fancy" page style to implement the FancyHdr headers

\newcommand{\HRule}{\rule{\linewidth}{0.5mm}} % New command to make the lines in the title page

% PDF meta-data
\hypersetup{pdftitle={\ttitle}}
% \hypersetup{pdfsubject=\subjectname}
\hypersetup{pdfauthor=\authornames}
% \hypersetup{pdfkeywords=\keywordnames}




% --------------------------------------------------





%----------------------------------------------------------------------------------------
%	TITLE PAGE
%----------------------------------------------------------------------------------------
\begin{comment}
\begin{titlepage}
\begin{center}

\HRule \\[0.4cm] % Horizontal line
{\huge \bfseries \ttitle}\\[0.4cm] % Thesis title
\HRule \\[1.5cm] % Horizontal line
 
\large \textit{A thesis submitted in fulfillment of the requirements\\ for the degree of \degreename}\\[0.3cm] % University requirement text
\textit{by}\\[0.4cm]

%\href{http://home.iitk.ac.in/~saiwal}{\authornames}

\vfill
\graphicspath{ {./Figures/} }
\begin{figure}[hb]
  \centering
  \includegraphics[width=0.4\linewidth]{Images/redlogo.png}
\end{figure}

\DEPTNAME\\ % Research group name and department name
\textsc{ \UNIVNAME}\\[1.5cm] % University name
\large \today\\[2cm] % Date


\end{center}

\end{titlepage}
\end{comment}
%\begin{document}
%\begin{titlepage}


%\include{title1}
%\end{titlepage}
%\myemptypage
%\afterpage{\null\newpage}
%\include{blank_page_1}
%\newpage
%\newpage
\include{title}

% \newpage
% \myemptypage
\frontmatter
\setstretch{1.3}
\pagestyle{plain}
% %----------------------------------------------------------------------------------------
% %	DEDICATION
% %----------------------------------------------------------------------------------------
% %
% \setstretch{1.3} % Return the line spacing back to 1.3
% \addtocontents{toc}{\vspace{2em}} % Add a gap in the Contents, for aesthetics
% %

% %
% \dedicatory{\textit{Dedicated to my beloved parents.}} % Dedication text
% % \thispagestyle{empty}
% %

% % \myemptypage

% %----------------------------------------------------------------------------------------
% Approval page
%----------------------------------------------------------------------------------------
\thispagestyle{empty} % Page style needs to be empty for this page

\addtotoc{Approval}
\addtocontents{toc}{}
% \pagenumbering{roman}
% \cfoot{\thepage}
% \setcounter{page}{3}
\vspace{-1in}
\begin{center}
{\LARGE  {\bf Dissertation Approval}}
\end{center}
\vspace*{0.1in} \noindent This dissertation entitled {\large \bf \dissertationTitle} by {\bf \large \authornames}, Roll No. \authorrollno, is approved for the degree of  {\large \bf \degree} from the Indian Institute of Technology Bombay.\\\\

\vspace{2cm}

  \hspace*{0.8cm}\ldots\ldots \ldots\ldots \ldots\ldots \ldots\ldots  \hspace{2in}  \ldots\ldots \ldots\ldots \ldots\ldots \ldots\ldots \\
  \hspace*{0.4cm}\hspace*{0cm} Prof. xxxxxxxxx xxxxxxx \hspace{1.85in} Prof. xxxxxxx xxxxxxxxx\\
  \hspace*{0.8cm}\hspace*{0.9cm} (Examiner 1) \hspace{2.8in} (Examiner 2)\\\\

\vspace{2cm}

  \hspace*{0.8cm}\ldots\ldots \ldots\ldots \ldots\ldots \ldots\ldots
  % \hspace{2in}  \ldots\ldots \ldots\ldots \ldots\ldots \ldots\ldots
  \\
  \hspace*{0.8cm}\hspace*{0cm} Prof. xxxxxx xxxxxxx 
  % \hspace{2.25in} Chairman name
  \\
  \hspace*{0.8cm}\hspace*{0.9cm} (Supervisor) 
  % \hspace{2.89in} (Chairman)
  \\\\

\vspace{1.5cm}

\noindent
 \hspace*{0.8cm}Defence Date: <Defence Date>\\\\
 \hspace*{0.8cm}Place: IIT Bombay


\ifnum \value{print}>0 {\myemptypage} \else {\clearpage} \fi

%-------------------------------------------------------------
%	DECLARATION PAGE
%	Your institution may give you a different text to place here
%----------------------------------------------------------------------------------------

% ------------------------------------------------
% Professor's declaration
\ifnum \value{print}>0 {\setcounter{page}{3}} \fi

\Declaration{\addtocontents{toc}{} % Add a gap in the Contents, for aesthetics
\vspace*{0.1in}
This is to certify that the dissertation entitled \textbf{``\dissertationTitle''}, submitted by \textbf{\authornames} to the Indian Institute of Technology Bombay, for the award of the degree of \textbf{\degree} in Electrical Engineering, is a record of the original, bona fide research work carried out by him under our supervision and guidance. The dissertation has reached the standards fulfilling the requirements of the regulations related to the award of the degree.

\par The results contained in this dissertation have not been submitted in part or in full to any other University or Institute for the award of any degree or diploma to the best of our knowledge.
\bigskip
\bigskip
\bigskip
\bigskip


\vspace*{2.5cm} %\vfill
%\begin{multicols}{2}
%\begin{center}
%\begin{flushleft}
%Date:
%\end{flushleft}
%\textbf{Prof. Akhilesh Kumar Maurya}\\
%Department of Civil Engineering,\\
%Indian Institute of Technology Guwahati.
%\end{center}
\begin{flushleft}
%\begin{flushleft}
%Date:
%\end{flushleft}
\ldots \ldots \ldots \ldots \ldots \ldots \ldots \ldots \ldots \ldots \ldots \ldots\\
\vspace{2mm}
\textbf{Prof. xxxxxx xxxxxxx}\\
Department of Electrical Engineering,\\
Indian Institute of Technology Bombay.
\end{flushleft}
%\end{multicols}
\vfill{}
}

\ifnum \value{print}>0 {\myemptypage} \else {\clearpage} \fi

% ------------------------------------------------
% Student declaration
\ifnum \value{print}>0 {\setcounter{page}{5}} \fi

\addtotoc{Declaration} % Add the "Scope for future work" page entry to the Contents

{\addtocontents{toc}{} % Add a gap in the Contents, for aesthetics

\begin{center}
    \textbf{\LARGE {Declaration}}
\end{center}}
\vspace*{0.1in}
I declare that this written submission represents my ideas in my own words. Where others' ideas and words have been included, I have adequately cited and referenced the original source. I declare that I have adhered to all principles of academic honesty and integrity and have not misrepresented or fabricated, or falsified any idea/data/fact/source in my submission. I understand that any violation of the above will cause disciplinary action by the Institute and can also evoke penal action from the source which has thus not been properly cited or from whom proper permission has not been taken when needed.

\bigskip
\bigskip
\bigskip
\bigskip


\vspace*{2.5cm} %\vfill
\begin{flushleft}
\ldots \ldots \ldots \ldots \ldots \ldots \ldots \ldots \ldots\\
\vspace{2mm}
\textbf{\authornames} \\
Roll No.: \authorrollno \\
Date: <Date>\\
Place: IIT Bombay
\end{flushleft}

\ifnum \value{print}>0 {\myemptypage} \else {\clearpage} \fi

% \setstretch{1.3} % Reset the line-spacing to 1.3 for body text (if it has changed)

%----------------------------------------------------------------------------------------
%	ABSTRACT PAGE
%----------------------------------------------------------------------------------------
\ifnum \value{print}>0 {\setcounter{page}{7}} \fi

\addtotoc{Abstract} % Add the "Abstract" page entry to the Contents 
\addtocontents{toc}{} % Add a gap in the Contents, for aesthetics
\thispagestyle{plain}
\abstract{
\setstretch{1.5}
    %%%%%%%%%%%%%%%%%
\vspace*{0.3in}
\lipsum[1-5]

\vspace*{2cm}

\textit{Index Terms} --- \lipsum[6]
    %%%%%%%%%%%%%%%%
}
\thispagestyle{plain}
\clearpage


% \myemptypage
%----------------------------------------------------------------------------------------
%	LIST OF CONTENTS/FIGURES/TABLES PAGES
%----------------------------------------------------------------------------------------

% \setcounter{page}{9}
\ifnum \value{print}=0 {\setcounter{page}{6}}
\pagestyle{fancy} % The page style headers have been "empty" all this time, now use the "fancy" headers as defined before to bring them back

\chead{\emph{Contents}} % Set the left side page header to "Contents"

\ifnum \value{print}>0 {\tableofcontents} \else {\setcounter{page}{6}\let\cleardoublepage\clearpage\tableofcontents } \fi

%combine figues and tables under tables, rename title appropriately
% \makeatletter
% \renewcommand*{\ext@figure}{lot}
% \let\c@figure\c@table
% \let\ftype@figure\ftype@table
% \let\listoftableandfigures\listoftables
% \renewcommand*\listtablename{List of Figures and Tables}
% \makeatother

\chead{\emph{List of Figures}} % Set the left side page header to "List of Figures"
\ifnum \value{print}>0 {\listoffigures} \else {\setcounter{page}{10}\let\cleardoublepage\clearpage\listoffigures} \fi
% \listoffigures % Write out the List of Figures

\chead{\emph{List of Tables}} % Set the left side page header to "List of Tables"
\ifnum \value{print}>0 {\listoftables} \else {\setcounter{page}{12}\let\cleardoublepage\clearpage\listoftables} \fi
% \listoftables % Write out the List of Tables

%----------------------------------------------------------------------------------------
%	ABBREVIATIONS
%----------------------------------------------------------------------------------------

\clearpage % Start a new page

\setstretch{1.5} % Set the line spacing to 1.5, this makes the following tables easier to read

\chead{\emph{Abbreviations}} % Set the left side page header to "Abbreviations"
\listofsymbols{ll} % Include a list of Abbreviations (a table of two columns)
{

% \textbf{SM} & \textbf{S}afety \textbf{M}argin\\

\textbf{CMOS}  & \textbf{C}omplementary \textbf{M}etal \textbf{O}xide \textbf{S}emiconductor Field Effect Transistors\\
\textbf{SOI}  & \textbf{S}ilicon-\textbf{o}n-\textbf{I}nsulator \\
\textbf{AI} & \textbf{A}rtificial \textbf{I}ntelligence \\
\textbf{ADC} & \textbf{A}nalog-to-\textbf{D}igital \textbf{C}onverter \\
\textbf{BTBT}  & \textbf{B}and-\textbf{T}o-\textbf{B}and \textbf{T}unnelling \\
\textbf{CTDSM} & \textbf{C}ontinuous-\textbf{T}ime \textbf{D}elta-\textbf{S}igma \textbf{M}odulator \\
\textbf{SNN} & \textbf{S}piking \textbf{N}eural \textbf{N}etwork \\
\textbf{LIF} & \textbf{L}eaky \textbf{I}ntegrate-and-\textbf{F}ire \\
\textbf{SFA} & \textbf{S}pike-\textbf{F}requency \textbf{A}daptation \\
\textbf{DSM} & \textbf{D}elta-\textbf{S}igma \textbf{M}odulator \\
\textbf{PMOS}  & \textbf{p}-type \textbf{MOS}FET \\
\textbf{NMOS}  & \textbf{n}-type \textbf{MOS}FET \\
\textbf{SCL} & \textbf{S}emi-\textbf{C}onductor \textbf{L}aboratory, Government of India\\
\textbf{LSM} & \textbf{L}iquid \textbf{S}tate \textbf{M}achine \\
\textbf{LF} & \textbf{L}oop \textbf{F}ilter \\
\textbf{RC} & \textbf{R}esistor-\textbf{C}apacitor \\
\textbf{UGF} & \textbf{U}nity-\textbf{G}ain \textbf{F}requency \\
\textbf{SR} & \textbf{S}lew \textbf{R}ate \\
\textbf{DAC} & \textbf{D}igital-to-\textbf{A}nalog \textbf{C}onverter \\
\textbf{NTF} & \textbf{N}oise \textbf{T}ransfer \textbf{F}unction \\
\textbf{GF45RFSOI} & \textbf{G}lobal\textbf{F}oundaries \textbf{45} nm \textbf{RF} \textbf{SOI} \\
\textbf{MOSFET} & \textbf{M}etal \textbf{O}xide \textbf{S}emiconductor \textbf{F}ield \textbf{E}ffect \textbf{T}ransistor \\
\textbf{DCR} & \textbf{D}uty-\textbf{C}ycled \textbf{R}esistor \\
\textbf{DC} & \textbf{D}irect \textbf{Current} \\
\textbf{RS} & \textbf{R}eset-\textbf{S}et \\
\textbf{FFT} & \textbf{F}ast \textbf{F}ourier \textbf{T}ransform \\
\textbf{SFDR} & \textbf{S}purious-\textbf{F}ree \textbf{D}ynamic \textbf{R}ange \\
\textbf{SNDR} & \textbf{S}ignal-to-\textbf{N}oise-and-\textbf{D}istortion \textbf{R}atio \\
\textbf{ENOB} & \textbf{E}ffective \textbf{N}umber-\textbf{o}f-\textbf{B}its \\
\textbf{OSR} & \textbf{O}ver\textbf{s}ampling \textbf{R}atio \\
\textbf{FoM} & \textbf{F}igure \textbf{o}f \textbf{M}erit \\
\textbf{DTDSM} & \textbf{D}iscrete-\textbf{T}ime \textbf{D}elta-\textbf{S}igma \textbf{M}odulator \\
\textbf{ERBW} & \textbf{E}ffective \textbf{R}esolution \textbf{B}and\textbf{w}idth \\
\textbf{BW} & \textbf{B}and\textbf{w}idth \\
\textbf{MOS} & \textbf{M}etal \textbf{O}xide \textbf{S}emiconductor \\
\textbf{SoC} & \textbf{S}ystem-\textbf{o}n-a-\textbf{C}hip \\
\textbf{ANN} & \textbf{A}rtificial \textbf{N}eural \textbf{N}etwork \\
\textbf{RNN} & \textbf{R}ecurrent \textbf{N}eural \textbf{N}etwork \\
\textbf{PseudoR} & \textbf{Pseduo}-\textbf{R}esistor \\
\textbf{PDK} & \textbf{P}rocess \textbf{D}esign \textbf{K}it \\
\textbf{IoT} & \textbf{I}nternet \textbf{o}f \textbf{T}hings \\
\textbf{HH} & \textbf{H}odgkin-\textbf{H}uxley \\
\textbf{AdEx} & \textbf{Ad}aptive \textbf{Ex}ponential Integrate-and-Fire \\
\textbf{VLSI} & \textbf{V}ery-\textbf{L}arge-\textbf{S}cale \textbf{I}ntegration \\
\textbf{MeLoDe} & \textbf{Me}mory, \textbf{Lo}gic Device and \textbf{De}sign \\
\textbf{ISCAS} & IEEE \textbf{I}nternational \textbf{S}ymposium on \textbf{C}ircuits \textbf{a}nd \textbf{S}ystems \\



}

%----------------------------------------------------------------------------------------

%----------------------------------------------------------------------------------------
%	SYMBOLS
%----------------------------------------------------------------------------------------

\clearpage % Start a new page
\setstretch{1.5} % Set the line spacing to 1.5, this makes the following tables easier to read

\chead{\emph{Symbols}} % Set the left side page header to "Symbols"

\listofnomenclature{lll} % Include a list of Symbols (a two column table)
{
$mV_{pp}$ & Millivolt peak-to-peak \\
$I_{in}$ & Input current \\
$V_{GS}$ & Gate-source voltage \\
$V_{th}$ & Threshold voltage \\
$I_{DS}$ & Drain-source current \\
$g_m$ & MOSFET Self-Transconductance \\
$V_{S}$ & Source voltage \\
$V_{G}$ & Gate voltage \\
$V_{D}$ & Drain voltage \\
$M_{(.)}$ & MOSFET labelled (.) \\
$FoM_{W}$ & Walden Figure-of-Merit \\
$FoM_{SNDR}$ & Schreier Figure-of-Merit \\
$FoM_{A}$ & Area * Walden Figure-of-Merit \\
$R_{(.)}$ & Resistor labelled (.) \\
$V_{tune}$ & Tuning terminal of a tunable Pseudo-resistor \\
$V_{control}$ & Control terminal for our proposed Pseudo-resistor neuron \\
$d\_tolXfet\_bc$ & Thick-oxide 472 nm length body contact device (X=n/p for NMOS/PMOS) \\
$C_{membrane}$ & Neuron membrane capacitance \\
$C_{sfa}$ & Spike-frequency adaptation capacitance \\



}



%----------------------------------------------------------------------------------------
%	THESIS CONTENT - CHAPTERS
%----------------------------------------------------------------------------------------

\mainmatter % Begin numeric (1,2,3...) page numbering

\pagestyle{fancy} % Return the page headers back to the "fancy" style

% Include the chapters of the thesis as separate files from the Chapters folder
% Uncomment the lines as you write the chapters

\input{Chapters/Chapter_1/Chapter1}


\part{My part1 label}
\label{part:part1}
\fancyhead[EC]{Part I. \emph{header label1}} % 
\input{Chapters/Chapter_2/Chapter2} 
\input{Chapters/Chapter_3/Chapter3}




\part[My part2 contents label]{My part2\\ label}
\label{part:part2}
\fancyhead[EC]{Part II. \emph{header label2}} % 
\input{Chapters/Chapter_4/Chapter4}



\clearpage % Start a new page
%----------------------------------------------------------------------------------------
%	THESIS CONTENT - APPENDICES
%----------------------------------------------------------------------------------------

% \addtocontents{toc}{\vspace{2em}} % Add a gap in the Contents, for aesthetics

% \appendix % Cue to tell LaTeX that the following 'chapters' are Appendices

% Include the appendices of the thesis as separate files from the Appendices folder
% Uncomment the lines as you write the Appendices

% \input{Appendices/AppendixA}
% \input{Appendices/AppendixB}
% \input{Appendices/AppendixC}
%\input{Appendices/AppendixD}

\addtocontents{toc}{\vspace{2em}} % Add a gap in the Contents, for aesthetics

\backmatter

%----------------------------------------------------------------------------------------
%	BIBLIOGRAPHY
%----------------------------------------------------------------------------------------

\label{References}

\chead{\emph{References}} % Change the page header to say "Bibliography"
%\usepackage{natbib}
%\renewcommand{\refname}{References}

\makeatletter
\interlinepenalty=10000     % don't split a single reference across page break

\renewcommand\bibname{References}
\bibliographystyle{ieeetr} % Use the "custom" BibTeX style for formatting the Bibliography
%\setcitestyle{authoryear,open={((},close={))}}

\bibliography{references} % The references (bibliography) information are stored in the file named "Bibliography.bib"

% \newcites{publ}{Publications by the candidate}

\clearpage
\section*{Publications by the candidate}
\begin{enumerate}
    \item A. Raut, A. Kadam, A. Singh, L. Somappa, M. S. Baghini and U. Ganguly, "A sub-100 nW Power, Compact CTDSM with a Band-to-Band Tunnelling Loop Filter," in \textit{2024 IEEE International Symposium on Circuits and Systems (ISCAS)} \{Presented in May 2024; Conference proceedings yet to be published\}
\end{enumerate}

\makeatother
% \newpage
% \include{LOP}
% \newpage
% \include{Biography}
\include{Acknowledgements}


\end{document}  


\end{document}